\documentclass[12pt]{extarticle}

% ===== Encodage & langue =====
\usepackage{fontspec}
\usepackage{polyglossia}
\setdefaultlanguage{french}

% ===== Police Verdana (compiler avec XeLaTeX ou LuaLaTeX) =====
\setsansfont{Verdana}
\renewcommand{\familydefault}{\sfdefault}

% ===== Interligne =====
\usepackage{setspace}
\setstretch{1.1}

% ===== Marges réduites =====
\usepackage{geometry}
\geometry{a4paper, top=1.8cm, bottom=1.8cm, left=2cm, right=2cm}

% ===== Mathématiques =====
\usepackage{amsmath}

% ===== Espacement réduit autour des équations =====
\setlength{\abovedisplayskip}{4pt}
\setlength{\belowdisplayskip}{4pt}
\setlength{\abovedisplayshortskip}{2pt}
\setlength{\belowdisplayshortskip}{2pt}

% ===== Espacement réduit des sections =====
\usepackage{titlesec}
\titlespacing*{\section}{0pt}{1.2ex}{0.6ex}

% ===== Titre du document =====
\title{\textbf{Correction du contrôle séquence 1 version 8}}
\author{}
\date{}

\begin{document}
\maketitle

% --------------------------------------------------
\section*{Exercice 1}
\begin{align*}
12 - 4 - 7 + 29 - 3 &= 8 - 7 + 29 - 3 \\
                     &= 1 + 29 - 3 \\
                     &= 30 - 3 \\
                     &= 27
\end{align*}

% --------------------------------------------------
\section*{Exercice 2}

\begin{align*}
\textbf{a)}
6 + 5 \times 7 &= 6 + 35 \\
               &= 41
\end{align*}

\textbf{b)}
\begin{align*}
56 - 9 \times 5 &= 56 - 45 \\
                &= 11
\end{align*}

\textbf{c)}
\begin{align*}
41 + 27 \div 9 &= 41 + 3 \\
               &= 44
\end{align*}

\textbf{d)}
\begin{align*}
6 + 8 \times 6 + 9 - 8 \div 2 &= 6 + 48 + 9 - 4 \\
                              &= 59
\end{align*}

% --------------------------------------------------
\section*{Exercice 3}
\begin{align*}
3 \times 8 + 4 \times 7 &= 24 + 28 \\
                        &= 52
\end{align*}

La dépense est de 52~€.

% --------------------------------------------------
\section*{Exercice 4}

\textbf{a)} $37 - (5 + 9) \div 2$
\begin{align*}
&= 37 - 14 \div 2 \\
&= 37 - 7 \\
&= 30
\end{align*}

\textbf{b)} $(5 + 8) \times 6 - 7 + 9$
\begin{align*}
&= 13 \times 6 - 7 + 9 \\
&= 78 - 7 + 9 \\
&= 80
\end{align*}

\textbf{c)} $(28 - (9 + 7)) \times 5$
\begin{align*}
&= (28 - 16) \times 5 \\
&= 12 \times 5 \\
&= 60
\end{align*}

\textbf{d)} $6 \times (29 - (8 + 15)) - 7$
\begin{align*}
&= 6 \times (29 - 23) - 7 \\
&= 6 \times 6 - 7 \\
&= 36 - 7 \\
&= 29
\end{align*}

% --------------------------------------------------
\section*{Exercice 5}
\begin{align*}
19 - 5 \times 3 &= 19 - 15 \\
                &= 4
\end{align*}

Une gomme coûte 4~€.

% --------------------------------------------------
\section*{Exercice bonus 1}

\textbf{a)} $7 \times 6 + 17$

\textbf{b)} $9 \times (8 + 5)$

% --------------------------------------------------
\section*{Exercice bonus 2}

\textbf{a)} La somme du produit de 8 par 6 et de 9.

\textbf{b)} Le produit de la somme de 7 et de 6 par 8.

\end{document}